% =======================================================
% 2. TERMINI
% =======================================================
\section{Termini}
 
% ------------- A -------------
\phantomsection
\addcontentsline{toc}{subsection}{A}
\subsection*{A}
\begin{itemize}[leftmargin=*]
      \item \textbf{Actual Cost (AC):} Costo effettivamente sostenuto per il lavoro 
      svolto entro una certa data.
      \item \textbf{Adapter:} Design pattern strutturale che permette a interfacce incompatibili di collaborare, traducendo l'interfaccia di una classe esistente in un'altra attesa dal client. Nel progetto è utilizzato per disaccoppiare il dominio applicativo dai dettagli di uno specifico provider LLM o libreria esterna.
      \item \textbf{Amministratore:} Ruolo del gruppo responsabile di assicurare
      l'efficienza di procedure, strumenti e tecnologie a supporto del
      \textit{way of working}; redige e mantiene aggiornate le Norme di
      Progetto e il Piano di Progetto; gestisce la creazione delle issue tramite Google Form; verifica,
      prima della chiusura di ogni sprint, che i tempi effettivi nel foglio di
      calcolo siano stati compilati correttamente.
      \item \textbf{Analisi dei Requisiti:} Documento essenziale che fissa i requisiti che il fornitore si impegna a soddisfare in accordo con il proponente. Traduce i bisogni dell'utente descrivendo esclusivamente \textit{cosa} il sistema deve fare, senza entrare in alcun dettaglio implementativo (il \textit{come}). Costituisce la base contrattuale e il documento fondamentale per l'intero ciclo di vita del progetto, aggiornato costantemente fino alla milestone di Product Baseline (PB).
      \item \textbf{Analisi Dei Rischi:} Sezione del Piano di Progetto che raccoglie l'elenco dei rischi identificati, il loro stato, la probabilità, l'impatto e le relative strategie di mitigazione e contingenza.
      \item \textbf{Analista:} Ruolo del gruppo responsabile delle attività di
      analisi dei requisiti. Attivo dalla baseline RTB.
      \item \textbf{Angular:} Framework frontend sviluppato da Google per la realizzazione di applicazioni web dinamiche e strutturate, basato su TypeScript.
      \item \textbf{API (Application Programming Interface):} Insieme di regole e strumenti che permettono a diversi sistemi software di comunicare tra loro. Nel contesto del progetto, consente l'invocazione di servizi esterni di LLM tramite chiamate strutturate.
      \item \textbf{API Key:} Chiave di autenticazione utilizzata per autorizzare l’accesso a servizi o API esterne.
      \item \textbf{Architecture Decision Record (ADR):} Documento sintetico utilizzato per registrare una decisione architetturale rilevante, includendo il contesto, le alternative valutate, la decisione presa e le relative motivazioni. Nel progetto è previsto per la fase PB, quando le scelte architetturali saranno consolidate.
      \item \textbf{Architettura a Layer:} Pattern architetturale che organizza il sistema in livelli sovrapposti (es. Presentazione, Business Logic), ciascuno dipendente solo dal livello sottostante. Adottata per il Backend a seguito di un confronto esplicito con l'azienda proponente (decisione VE-8.1), in alternativa all'architettura Esagonale.
      \item \textbf{Architettura Esagonale (Ports and Adapters):} Pattern architetturale che isola la logica di dominio dai dettagli tecnologici esterni (framework, database, API) tramite porte e adapter, favorendo testabilità e sostituibilità dei componenti. Valutata come alternativa all'architettura a Layer per il Backend e scartata per la scala contenuta del progetto (decisione VE-8.1).
      \item \textbf{Artefatto:} Prodotto concreto generato durante lo svolgimento del progetto, come documento, diagramma, codice sorgente, script, file PDF, verbale o materiale di presentazione.
      \item \textbf{Appunti (Clipboard):} Area temporanea di memoria del sistema operativo utilizzata per conservare dati copiati o tagliati dall'utente, rendendoli disponibili per operazioni di incolla. Nel progetto è utilizzata per trasferire testo da e verso l'editor.
      \item \textbf{Asincrono:} Modello di esecuzione in cui operazioni e processi possono essere svolti senza bloccare il flusso principale del programma.   
      \item \textbf{Attore:} In ambito UML, rappresenta un'entità esterna (un utente umano o un altro sistema) che interagisce direttamente con il sistema software per raggiungere uno specifico obiettivo.
      \item \textbf{AWS (Amazon Web Services):} Piattaforma di cloud computing che fornisce servizi on-demand come elaborazione, storage e database.
      \end{itemize}
 
% ------------- B -------------
\phantomsection
\addcontentsline{toc}{subsection}{B}
\subsection*{B}
\begin{itemize}[leftmargin=*]
      \item \textbf{Budget at Completion (BAC):} Budget totale pianificato per 
      l'intero progetto.
      \item \textbf{Budget Variance (BV):} Scostamento tra il costo consuntivato e quello preventivato per il lavoro svolto entro una certa data; sinonimo di Cost Variance (CV), utilizzato in alcuni documenti di fase RTB.
      \item \textbf{Backlog:} Lista ordinata per priorità di tutte le attività, funzionalità
      e requisiti da realizzare nel corso del progetto. Nel contesto del gruppo viene
      gestita tramite GitHub Project e aggiornata ad ogni Sprint
      Planning.
      \item \textbf{Baseline:} Versione stabile e approvata di un insieme di
          documenti o artefatti, fissata in un determinato momento del progetto
          e utilizzata come riferimento ufficiale per le fasi successive. Nel
          contesto del progetto didattico le baseline principali sono: candidatura
          (CC), Requirements and Technology Baseline (RTB) e Product Baseline (PB).
      \item \textbf{Backend:} Parte dell'applicazione che gestisce logica server, API e
      servizi di supporto non direttamente visibili all'utente.
      \item \textbf{Brainstorming:} Tecnica di generazione collettiva delle idee adottata
      dal gruppo come metodo decisionale principale. Si articola in tre fasi:
      raccolta libera delle proposte senza valutazione immediata, confronto
      strutturato con analisi di vantaggi e svantaggi, scelta collegiale della
      soluzione ritenuta più adeguata con registrazione nel verbale.
      \item \textbf{Branch:} Linea di sviluppo indipendente all'interno di un
          repository Git, utilizzata per lavorare su modifiche isolate senza
          influenzare il ramo principale (\texttt{main}).
      \item \textbf{Branch Protection Rules:} Regole di protezione configurate su GitHub per limitare le modifiche dirette a un branch. Nel progetto sono utilizzate sul branch \texttt{main} per richiedere almeno una review positiva, bloccare i force push e restringere l'eliminazione del branch protetto.
      \item \textbf{Browser Compatibility:} Metrica di prodotto che misura il numero di
browser target su cui il sistema è verificato o testato.
      \item \textbf{Budget:} Quantità di risorse economiche disponibili per il progetto o per uno specifico sprint, calcolata sulla base delle ore preventivate e dei costi orari associati ai ruoli.
      \item \textbf{Bug:} Comportamento errato o non conforme rispetto a quanto atteso
      per una funzionalità o componente software.
      \item \textbf{Builder:} Design pattern creazionale che separa la costruzione incrementale di un oggetto complesso dalla sua rappresentazione finale, tramite metodi fluenti che restituiscono il builder stesso. Nel progetto è utilizzato per assemblare in modo incrementale il \texttt{Prompt} da inviare all'LLM.
      \item \textbf{Burndown Chart:} Grafico che mostra l'andamento del lavoro residuo nel
      corso di uno sprint o dell'intero progetto, confrontando il ritmo di avanzamento
      effettivo con quello pianificato. Utilizzata per monitorare il rischio di ritardo
      rispetto alle milestone.
\end{itemize}
 
% ------------- C -------------
\phantomsection
\addcontentsline{toc}{subsection}{C}
\subsection*{C}
\begin{itemize}[leftmargin=*]
      \item \textbf{Candidatura (CC):} Fase iniziale del progetto didattico, antecedente all'aggiudicazione del capitolato, in cui il gruppo si propone come fornitore producendo la Valutazione dei Capitolati, la Lettera di Presentazione e il Preventivo Costi e Assunzione Impegni. Non è un documento, ma una baseline di fase.      
      \item \textbf{Cerimonie:} Insieme degli incontri strutturati previsti dal
      framework Scrum: Sprint Planning, Sprint Review, Sprint
      Retrospective e Daily Stand-up. Nel gruppo le prime tre si svolgono il martedì in sequenza fissa; il Daily Stand-up è gestito in modo asincrono.
      \item \textbf{Change Request:} Richiesta di cambiamento relativa a documenti, requisiti, tecnologie o processi. Deve essere valutata rispetto al suo impatto, approvata dal ruolo competente, tracciata tramite issue e verificata tramite Pull Request.
      \item \textbf{Chatbot:} Sistema software che consente l'interazione con l'utente tramite linguaggio naturale. Nel progetto può fungere da mediatore tra le richieste dell'utente e le funzionalità disponibili, traducendo intenzioni espresse liberamente in azioni eseguibili.
      \item \textbf{Casi d'uso:} Tecnica di modellazione che descrive le funzionalità del sistema tramite una "visione esterna" (come percepite dall'attore), ignorandone l'implementazione. A livello tecnico è definito come un insieme di scenari (sequenze di azioni che includono uno scenario principale e possibili scenari alternativi o di errore) che condividono lo stesso scopo finale per l'utente.
      \item  \textbf{CI/CD (Continuous Integration / Continuous Deployment):} Pratica di sviluppo software che prevede l’integrazione frequente del codice in un repository condiviso e la sua automazione di build, test e distribuzione.
      \item \textbf{Ciclo Di Vita Documentale:} Sequenza di stati attraversati da un documento dalla creazione al rilascio. Nel gruppo si articola in: In Stesura, In Verifica, Verificato, Approvato e Rilasciato.      \item \textbf{Code Coverage:} Metrica che misura la percentuale di codice sorgente eseguita durante l’esecuzione dei test automatici.
      \item \textbf{CodeMirror:} Libreria open source per l'implementazione di editor di testo all'interno del browser, utilizzata come componente centrale dell'editor Markdown del Frontend, integrata da un insieme di pacchetti ufficiali per le funzionalità specifiche richieste (sintassi Markdown, tema, gestione dello stato).
      \item \textbf{Client-Server:} Pattern architetturale che separa un'applicazione in due componenti distinti -- il client, che gestisce l'interfaccia utente, e il server, che espone servizi e logica applicativa -- comunicanti tramite un protocollo di rete. Adottato come architettura generale del sistema per separare la gestione dell'editor (Frontend) dall'orchestrazione delle chiamate LLM (Backend).
      \item \textbf{Colli di bottiglia:} Fase o componente di un processo che rallenta l'efficienza complessiva, limitando la capacità operativa del team o del sistema.
      \item \textbf{Command:} Design pattern comportamentale che incapsula una richiesta come oggetto, permettendo di parametrizzare, accodare o annullare operazioni in modo uniforme. Nel progetto è utilizzato per implementare Undo/Redo delle modifiche testuali, incluse quelle generate dall'AI.
      \item \textbf{Commit:} Operazione Git che registra un insieme di modifiche nel
      repository locale, associandovi un messaggio descrittivo. Nel flusso del gruppo
      ogni commit deve seguire la convenzione
      \texttt{<tipo>(<posizione>): <messaggio> (\#id)}, così da garantire
      tracciabilità rispetto alla issue collegata.
      \item \textbf{Comment Density:} Metrica di qualità del prodotto che misura il
      rapporto tra righe di commento e righe di codice totali.
      \item \textbf{Compito:} Obiettivo espresso dall'utente in linguaggio naturale (es.\ ``applicare il cappello critico''), che il sistema interpreta per determinare la funzionalità da eseguire.
      \item \textbf{Complessità Ciclomatica:} Metrica di qualità del software che misura
      la complessità dei flussi di controllo di una funzione o di un modulo, utile per
      valutare manutenibilità e difficoltà di verifica.
      \item \textbf{Configurazione:} Insieme degli elementi che compongono lo stato controllato del progetto, inclusi documenti, codice, template, asset, script, file PDF e configurazioni del repository.
      \item \textbf{Consuntivo:} Documento o sezione del Piano di Progetto che riporta, al termine di ogni sprint, le ore effettivamente
      impiegate da ciascun membro per ciascun ruolo, il relativo costo
      effettivo e lo scostamento rispetto al preventivo.
      \item \textbf{Container:} Unità software isolata che include codice, dipendenze e configurazioni necessarie all’esecuzione di uno specifico servizio.
      \item \textbf{Containerizzazione:} Processo di distribuzione ed esecuzione del software tramite container isolati e portabili.
      \item \textbf{Conventional Commits:} Convenzione per la scrittura dei messaggi di
      commit che associa a ogni modifica un tipo, uno scope e una descrizione
      sintetica. Nel progetto viene adottata in forma adattata al workflow del gruppo.
      \item \textbf{Cruscotto Di Qualità:} Strumento di raccolta e visualizzazione delle metriche di qualità di processo e di prodotto, utilizzato per monitorare il raggiungimento degli obiettivi definiti nel Piano di Qualifica.
      \item \textbf{Contingenza:} Insieme di azioni correttive attivate quando un rischio si manifesta effettivamente, con l'obiettivo di ridurne gli effetti sul progetto.
      \item \textbf{Cost Performance Index (CPI):} Indice di performance dei costi, se >1 
      si è sotto budget, se <1 si è sopra budget.
      \item \textbf{Cost Variance (CV):} Varianza dei costi. 
      Differenza tra EV e AC. Valore positivo indica risparmio, 
      negativo indica superamento del budget.
\end{itemize}
 
% ------------- D -------------
\phantomsection
\addcontentsline{toc}{subsection}{D}
\subsection*{D}
\begin{itemize}[leftmargin=*]
      \item \textbf{Daily Stand-up:}
      Momento di coordinamento rapido tra i membri del gruppo utilizzato durante
      lo sprint per condividere eventuali avanzamenti rilevanti, impedimenti,
      dubbi o richieste di supporto. Nel gruppo viene svolto prevalentemente in
      modalità asincrona tramite WhatsApp. Lo stato di avanzamento delle attività
      è tracciato principalmente tramite GitHub Projects, issue e pull request.
      \item \textbf{Database Relazionale:} Sistema di gestione dei dati organizzato in tabelle correlate tra loro tramite relazioni strutturate.
      \item \textbf{Debito Tecnico:} Insieme di scelte progettuali o implementative
      sub-ottimali accumulate durante lo sviluppo per ragioni di tempo o priorità,
      che richiedono interventi correttivi futuri (refactoring) e che, se non gestite,
      rallentano la velocity degli sprint successivi.
      \item \textbf{Decision Tree (Albero Di Decisione):} Modello strutturato ad albero utilizzato per rappresentare decisioni e le loro possibili conseguenze.
      \item \textbf{Definition of Done (DoD):} Lista esplicita di criteri che un task
      deve soddisfare per essere considerato completato a tutti gli effetti. Nella fase
      RTB ha incluso criteri documentali e criteri minimi per il codice del Proof of
      Concept; nella fase PB è stata estesa a test automatici e copertura del
      codice.
      \item \textbf{Dependency Injection (DI):} Pattern creazionale che implementa l'Inversion of Control fornendo le dipendenze di un componente dall'esterno (tipicamente tramite costruttore), invece di farle creare al componente stesso. Nel progetto è utilizzato sia lato Backend (tramite il meccanismo nativo di FastAPI) sia lato Frontend, dove è stato introdotto autonomamente dal gruppo in assenza di supporto nativo di Vue.js.
      \item \textbf{Deployment:} Processo di rilascio e distribuzione di un’applicazione in un ambiente di esecuzione.
      \item \textbf{Deployment Automatizzato:} Distribuzione di un'applicazione eseguita senza intervento manuale, tramite strumenti di orchestrazione e configurazione predefinita. Nel progetto è realizzato tramite Docker Compose, che avvia ed orchestra i container Frontend e Backend a partire da una configurazione dichiarativa.
      \item \textbf{Diario Di Bordo:} Momento di rendicontazione periodica al committente
      in cui il gruppo espone il rapporto tra costi sostenuti e avanzamento
      conseguito, comunicando le difficoltà riscontrate. Si svolge secondo il
      calendario stabilito dai docenti.
      \item \textbf{Discord:} Piattaforma di comunicazione vocale e testuale utilizzata
            dal gruppo come ambiente principale per le riunioni interne e il coordinamento
            quotidiano.
      \item \textbf{Distant Writing:} Funzionalità di generazione automatica del testo tramite prompt, in cui l’LLM produce contenuti estesi a partire da istruzioni concise fornite dall’utente.
      \item \textbf{Docker:} Piattaforma per creare, distribuire ed eseguire applicazioni
      tramite container.
      \item \textbf{Docker Compose:} Strumento per definire ed eseguire applicazioni
      composte da più container tramite un file di configurazione.
      \item \textbf{DTO (Data Transfer Object):} Oggetto utilizzato per trasferire dati strutturati tra diversi livelli o componenti di un sistema, tipicamente privo di logica applicativa propria. Nel progetto i DTO validano e normalizzano i payload delle richieste/risposte API.
\end{itemize}
 
% ------------- E -------------
\phantomsection
\addcontentsline{toc}{subsection}{E}
\subsection*{E}
\begin{itemize}[leftmargin=*]
      \item \textbf{Endpoint:} Punto di accesso esposto da un servizio web o da un'API, identificato da un URL specifico, tramite il quale un client può inviare richieste e ottenere risposte. Nel progetto indica l'endpoint API del servizio AI utilizzato per l'elaborazione delle richieste.
      \item \textbf{Embedding:} Rappresentazione vettoriale di informazioni testuali o semantiche utilizzata nei modelli di machine learning e negli LLM.
      \item \textbf{Estimate at Completion (EAC):} Stima del costo totale del 
      progetto al termine, basata sull'andamento attuale.
      \item \textbf{Earned Value (EV):} Valore guadagnato. 
      Budget del lavoro effettivamente completato entro una certa data. 
      Si calcola come somma dei costi previsti dei task completati
\end{itemize}
 
% ------------- F -------------
\phantomsection
\addcontentsline{toc}{subsection}{F}
\subsection*{F}
\begin{itemize}[leftmargin=*]
      \item \textbf{Fallback:} Strategia che prevede un comportamento o percorso alternativo da attivare quando l'operazione principale non va a buon fine. Nel progetto è utilizzato per la gestione degli errori di comunicazione con il servizio LLM esterno.
      \item \textbf{FastAPI:} Framework moderno Python per la creazione di API REST ad alte prestazioni, con supporto asincrono e generazione automatica della documentazione.
      \item \textbf{Feature:} Nuova funzionalità del prodotto, del Proof of Concept o di
      una sua componente.
      \item \textbf{File System:} Componente del sistema operativo responsabile della gestione dei file locali. Nel progetto può fungere da attore secondario per operazioni di upload e download delle note.
      \item \textbf{Foglio Di Calcolo:} Strumento condiviso del gruppo utilizzato per
      il monitoraggio aggregato delle misurazioni. Raccoglie i dati delle task
      registrati tramite Google Form e calcola automaticamente, per ogni sprint,
      preventivo e consuntivo di ore e costi per membro e per ruolo, oltre alle
      risorse residue. Le formule applicate sono: $\text{ore} = \text{minuti} / 60$
      e $\text{costo} = \text{ore} \times \text{costo orario del ruolo}$.
      \item \textbf{Force Push:} Operazione Git che riscrive la cronologia di un branch remoto. Nel progetto è bloccata sul branch \texttt{main} tramite branch protection rules.
      \item \textbf{Formattazione:} Insieme delle funzionalità dell'editor che
            consentono di applicare stili e struttura al testo in Markdown.
            Nel progetto include: grassetto, corsivo, sottolineato, barrato,
            titoli, link ipertestuali, elenchi puntati e numerati, tabelle e
            citazioni. La gestione delle immagini è esplicitamente esclusa dai
            requisiti.
      \item \textbf{Framework:} Struttura software di base o piattaforma che fornisce librerie e strumenti per facilitare e standardizzare lo sviluppo di applicazioni.
      \item \textbf{Frontend:} Parte dell'applicazione con cui l'utente interagisce
      direttamente tramite interfaccia grafica.
\end{itemize}
 
% ------------- G -------------
\phantomsection
\addcontentsline{toc}{subsection}{G}
\subsection*{G}
\begin{itemize}[leftmargin=*]
      \item \textbf{GitHub Actions:} Servizio di automazione integrato in GitHub. Nel progetto è utilizzato per automatizzare la creazione delle issue a partire dai dati raccolti tramite Google Form.
      \item \textbf{GitHub Pages:} Servizio di GitHub che consente di pubblicare siti
            web statici direttamente da un repository. Utilizzato dal gruppo per
            esporre la documentazione prodotta come sito web accessibile pubblicamente.
      \item \textbf{GitHub Project:} Strumento di gestione dei task integrato in
            GitHub, utilizzato per pianificare, tracciare e monitorare le attività
            del progetto tramite board e issue.
      \item \textbf{Google Calendar:} Servizio di calendario condiviso di Google,
            utilizzato dal gruppo per la gestione delle disponibilità dei membri e
            la pianificazione delle riunioni ricorrenti e straordinarie.
      \item \textbf{Google Form:} Strumento di Google utilizzato per la creazione
          guidata dei task; genera automaticamente una issue su GitHub e registra
          i dati nel foglio di calcolo condiviso per il monitoraggio delle
          misurazioni.
\end{itemize}
 
% ------------- H -------------
\phantomsection
\addcontentsline{toc}{subsection}{H}
\subsection*{H}
\begin{itemize}[leftmargin=*]
      \item \textbf{HTTPS:} Protocollo di comunicazione sicuro che utilizza cifratura TLS per proteggere i dati trasmessi in rete.
      \item \textbf{HTML:} (HyperText Markup Language) Linguaggio di formattazione utilizzato per strutturare e visualizzare documenti sul World Wide Web.
\end{itemize}
 
% ------------- I -------------
\phantomsection
\addcontentsline{toc}{subsection}{I}
\subsection*{I}
\begin{itemize}[leftmargin=*]
      \item \textbf{Indice di Gulpease}: Indice di leggibilità della lingua italiana che
      misura la facilità di lettura di un testo sulla base della lunghezza delle parole
      e delle frasi.
      \item \textbf{In Review:} Stato di una issue in GitHub Projects che indica che l'attività è stata completata dal Redattore ed è in attesa di verifica tramite Pull Request.
      \item \textbf{Integrazione Continua:} Pratica di sviluppo che prevede l'integrazione frequente delle modifiche nel repository principale, accompagnata da controlli automatici. Nel progetto è pianificata per la fase PB.
      \item \textbf{Interfaccia Utente (UI):} Insieme degli elementi grafici e interattivi attraverso cui l'utente interagisce con il sistema. Deve garantire semplicità, chiarezza e accessibilità.
      \item \textbf{ISO/IEC 12207:1995:} Standard internazionale
            (\textit{Information Technology~-- Software Life Cycle Processes})
            che definisce un framework di processi per il ciclo di vita del software,
            suddivisi in processi primari (es.\ fornitura, sviluppo), di supporto
            (es.\ documentazione, verifica, gestione della configurazione) e
            organizzativi (es.\ gestione, infrastruttura). Fornisce un quadro
            metodologico per strutturare, gestire e controllare le attività di
            sviluppo e manutenzione del software, applicabile tramite tailoring al
            contesto specifico del progetto.
      \item \textbf{Issue:} Unità di lavoro o segnalazione registrata su GitHub,
            utilizzata per tracciare attività, problemi o richieste. Nel flusso del
            gruppo ogni issue corrisponde a un task e viene generata automaticamente
            tramite Google Form.
      \item \textbf{Issue Tracking:} Processo di monitoraggio delle attività tramite issue GitHub, includendo stati progressivi, assegnatari, tempi previsti ed effettivi, e collegamento con il sistema di misurazione ore.
\end{itemize}
 
% ------------- J -------------
\phantomsection
\addcontentsline{toc}{subsection}{J}
\subsection*{J}
\begin{itemize}[leftmargin=*]
      \item \textbf{JSON:} (JavaScript Object Notation) Formato leggero per lo scambio di dati, facilmente leggibile dalle macchine e dagli esseri umani.
      \item \textbf{Just-in-time:} Principio organizzativo secondo cui ogni attività
          viene normata immediatamente prima di essere attuata, evitando
          pianificazioni eccessive in anticipo. Applicato al \textit{way of working},
          implica che le norme vengano estese e aggiornate in modo incrementale,
          a intervalli brevi, man mano che nuove attività vengono introdotte.
      \item \textbf{JWT:} Standard aperto per la trasmissione sicura di informazioni tra parti tramite token firmati digitalmente.
\end{itemize}
 
% ------------- K -------------
\phantomsection
\addcontentsline{toc}{subsection}{K}
\subsection*{K}
\textit{Nessun termine.}
 
% ------------- L -------------
\phantomsection
\addcontentsline{toc}{subsection}{L}
\subsection*{L}
\begin{itemize}[leftmargin=*]
      \item \textbf{Latenza:} Tempo necessario affinché un sistema risponda a una richiesta effettuata da un utente o da un servizio.
      \item \textbf{LLM (Large Language Model):} Modello avanzato di intelligenza
          artificiale addestrato su grandi quantità di testo per la comprensione,
          elaborazione e generazione di linguaggio umano in modo contestuale.
      \item \textbf{Lettera di presentazione:} Documento formale redatto dal gruppo per accompagnare la candidatura, contenente la motivazione della scelta del capitolato e i collegamenti ai documenti prodotti.
      \item \textbf{Logout:} Operazione con cui un utente termina volontariamente la propria sessione all'interno dell'applicazione, invalidando eventuali token o credenziali attive e impedendo ulteriori interazioni fino a una nuova autenticazione.
      \item \textbf{Lucidspark:} Strumento online per la creazione collaborativa di diagrammi, utilizzato dal gruppo per la modellazione UML e la rappresentazione dei Casi d'Uso.
\end{itemize}
 
% ------------- M -------------
\phantomsection
\addcontentsline{toc}{subsection}{M}
\subsection*{M}
\begin{itemize}[leftmargin=*]
      \item \textbf{Main:} Ramo (branch) principale del repository Git, che contiene esclusivamente le versioni stabili, verificate e approvate dei documenti e del codice.
      \item \textbf{Markdown:} Linguaggio di markup leggero che permette di formattare il testo utilizzando una sintassi semplice e intuitiva in puro testo semplice.
      \item \textbf{Matrice Dei Rischi:} Strumento grafico che classifica i rischi in base alla loro probabilità di occorrenza e al loro impatto sul progetto, utilizzato per prioritizzare le azioni di gestione.
      \item \textbf{Merge:} Operazione Git che integra le modifiche di un branch
      di lavoro nel branch principale (\texttt{main}). Nel flusso del gruppo
      il merge viene eseguito dall'Assegnatario del task (non dal Verificatore)
      dopo che la pull request ha ricevuto l'approvazione del Verificatore.
      \item \textbf{Merge Commit:} Modalità di integrazione di una Pull Request che mantiene la cronologia dei commit del branch di lavoro e aggiunge un commit di merge nel branch di destinazione.
      \item \textbf{Messaggi Precompilati:} Elementi interattivi (es.\ bottoni) che rappresentano azioni o richieste già definite, proposte all'utente per guidare l'interazione e ridurre l'ambiguità del linguaggio naturale.
      \item \textbf{Middleware:} Componente software intermedio che gestisce elaborazioni, sicurezza o controlli tra client e backend.
      \item \textbf{Milestone:} Punto di controllo formale nel calendario di progetto che
      segna il completamento di una fase o il raggiungimento di un obiettivo
      significativo. Nel progetto le milestone principali sono la consegna RTB
      (giugno--luglio 2026) e la consegna PB (agosto 2026).
      \item \textbf{Misurazione:} Raccolta strutturata di dati quantitativi relativi
          allo svolgimento delle attività del gruppo, comprensiva di: ruolo
          ricoperto, tempo previsto e tempo effettivo impiegato. I dati sono
          registrati nel foglio di calcolo condiviso e utilizzati per il
          monitoraggio dell'avanzamento.
      \item \textbf{Mitigazione:} Insieme di azioni preventive adottate per ridurre la probabilità di occorrenza o l'impatto di un rischio identificato.
      \item \textbf{Model-View-Controller (MVC):} Pattern architetturale che separa la logica di dominio (Model), la presentazione (View) e la gestione degli input utente (Controller) in tre componenti distinti e disaccoppiati. Adottato per il Frontend nella variante \textit{Pull Model}, in cui la View espone il proprio stato ed è il Controller a interrogarlo, anziché riceverlo push dal Model.
      \item \textbf{MongoDB:} Database NoSQL orientato ai documenti, open source,progettato per scalabilità orizzontale e alte prestazioni in contesti distribuiti.
      \item \textbf{MVP (Minimum Viable Product):} Versione del prodotto dotata delle
          funzionalità minime sufficienti per essere valutata dal proponente.
          Consente di verificare la bontà della visione iniziale, di raccogliere
          feedback fondato e di prendere decisioni informate per il completamento
          definitivo del prodotto.
\end{itemize}
 
% ------------- N -------------
\phantomsection
\addcontentsline{toc}{subsection}{N}
\subsection*{N}
\begin{itemize}[leftmargin=*]
      \item \textbf{Non Conformità:} Deviazione rispetto alle Norme di Progetto, al template, al Piano di Qualifica o agli accordi formalizzati nei verbali. Può essere bloccante o non bloccante.
      \item \textbf{Norme Di Progetto:} Documento gestionale interno che definisce
          le procedure operative, le convenzioni redazionali e gli strumenti adottati
          dal gruppo, organizzati secondo i processi del ciclo di vita definiti da
          ISO/IEC 12207:1995. Redatto e mantenuto dall'Amministratore; soggetto
          ad aggiornamento incrementale secondo il principio just-in-time.
      \item \textbf{NoSQL:} Categoria di database non relazionali progettati per garantire maggiore flessibilità e scalabilità nella gestione dei dati.
\end{itemize}
 
% ------------- O -------------
\phantomsection
\addcontentsline{toc}{subsection}{O}
\subsection*{O}
\begin{itemize}[leftmargin=*]
    \item \textbf{Observer:} Design pattern comportamentale che permette a un oggetto (Subject) di notificare automaticamente i propri cambiamenti di stato a una lista di osservatori (Observer) senza conoscerne il tipo concreto. Nel progetto realizza la propagazione degli aggiornamenti tra Model, View e Controller nell'architettura MVC del Frontend.
    \item \textbf{OCR (Optical Character Recognition):} Tecnologia che permette di riconoscere e convertire testo contenuto in immagini o documenti scansionati in dati testuali modificabili.
    \item \textbf{OdG Preliminare:} Acronimo di ``Ordine del Giorno Preliminare''.
          Indica l'elenco dei punti da trattare nella riunione, pubblicato dal
          Responsabile prima dell'inizio della seduta su un documento Google
          condiviso. Funge anche da base per la successiva redazione del verbale
          interno, in quanto nel medesimo documento vengono annotate le decisioni
          adottate durante la riunione.
    \item \textbf{Open Closed Principle (OCP):} Principio di progettazione orientata agli oggetti secondo cui un'entità software dovrebbe essere aperta all'estensione ma chiusa alla modifica. Nel progetto guida la scelta di pattern come Strategy e Simple Factory, che permettono di aggiungere nuove operazioni AI senza modificare il codice esistente.
    \item \textbf{OWASP (Open Web Application Security Project):} Fondazione no-profit che produce standard, strumenti e linee guida per la sicurezza delle applicazioni web.
\end{itemize}
 
% ------------- P -------------
\phantomsection
\addcontentsline{toc}{subsection}{P}
\subsection*{P}
\begin{itemize}[leftmargin=*]
      \item \textbf{Pair Programming:} Tecnica di sviluppo software in cui due membri del
      team lavorano insieme sullo stesso task, condividendo lo stesso
      strumento di lavoro: uno scrive (driver) e l'altro revisiona in tempo reale
      (navigator). Utilizzata come misura di mitigazione per i rischi legati
      all'inesperienza tecnologica.
      \item \textbf{Pannello contestuale:} Elemento dell'interfaccia utente che compare in risposta a un'azione dell'utente per mostrare informazioni, suggerimenti o funzionalità rilevanti rispetto al contesto corrente. Nel progetto viene utilizzato per presentare i risultati delle analisi AI o opzioni operative senza interrompere il flusso di lavoro.
      \item \textbf{PB (Product Baseline):} Seconda revisione formale di avanzamento
      del progetto, successiva alla RTB. Fissa il prodotto software completo,
      includendo architettura, codice e documentazione tecnica finale. Nel
      calendario del gruppo è prevista per agosto 2026.
      \item \textbf{PDCA (Plan-Do-Check-Act):} Ciclo di miglioramento continuo composto da pianificazione, esecuzione, controllo e azione correttiva. Nel progetto è utilizzato come riferimento per l'accertamento della qualità e il miglioramento del way of working.
      \item \textbf{Pedice G:} Convenzione tipografica adottata nelle norme di progetto
          per segnalare che un termine è definito nel Glossario. Viene applicato
          alla prima occorrenza del termine in ciascun documento tramite lo script
          Python di automazione (\texttt{python\_glossary\_automation/}).
          \textit{Esempio}: Termine\textsubscript{G}.
      \item \textbf{PII:} Informazioni personali identificabili che consentono l’identificazione diretta o indiretta di un individuo.
      \item \textbf{Piano di Progetto:} Documento gestionale e strategico mirato all'efficienza del team. Definisce l'organigramma, l'assegnazione e la rotazione dei ruoli, il preventivo dell'impegno orario e il consuntivo di periodo. Gestisce l'analisi dei rischi e struttura le attività in una pianificazione macroscopica a lungo termine (fissando le milestone a ritroso) e una pianificazione dettagliata a breve termine (guardando in avanti per i singoli sprint).
      \item \textbf{Piano Di Qualifica:} Documento gestionale esterno che definisce
      gli obiettivi quantitativi di qualità di processo e di prodotto, le
      metriche adottate per misurarli e il cruscotto di valutazione utilizzato
      per monitorarne il raggiungimento nel tempo. Include le retrospettive
      e le iniziative di automiglioramento del gruppo.
      \item \textbf{Planned Value (PV):} Valore pianificato. 
      Budget del lavoro che avrebbe dovuto essere completato entro una certa data 
      secondo la pianificazione di base.
      \item \textbf{PoC (Proof of Concept):} Artefatto eseguibile, da considerarsi
          usa-e-getta sul piano architetturale, realizzato all'inizio del progetto
          per valutare la fattibilità tecnologica del prodotto atteso rispetto a
          specifici sottoinsiemi di funzionalità individuati con il proponente.
          Costituisce la technology baseline e aiuta a consolidare l'analisi dei
          requisiti dal lato della soluzione.
      \item \textbf{Pop-up:} Finestra grafica temporanea che appare in sovrapposizione all'interfaccia principale per mostrare informazioni o risultati contestuali senza interrompere il flusso di lavoro dell'utente.
      \item \textbf{Portabilità:} Capacità di un software di essere eseguito su ambienti differenti senza modifiche sostanziali.
      \item \textbf{Postgres (PostgreSQL):} Sistema avanzato di gestione di database
          relazionali e ad oggetti, open source. Utilizzato come tecnologia di
          persistenza dei dati nel capitolato di riferimento.
      \item \textbf{Preventivo:} Stima delle ore e dei costi per ogni membro e ruolo
            per uno sprint, calcolata a partire dai tempi previsti delle task assegnate.
            Distinto dal \textit{Preventivo Dei Costi} del Preventivo dei Costi e Assunzione degli Impegni,
            che è il preventivo globale di progetto; il preventivo di sprint è prodotto
            automaticamente dal foglio di calcolo condiviso.
      \item \textbf{Preventivo Dei Costi:} Stima dell'impegno economico complessivo
            del progetto, calcolata moltiplicando le ore previste per ciascun ruolo
            per il relativo costo orario ufficiale. Presentato nel Preventivo dei Costi e Assunzione degli Impegni; il valore accettato all'aggiudicazione diventa limite
            superiore invalicabile durante lo svolgimento del progetto.
      \item \textbf{Progettista:} Ruolo del gruppo responsabile delle attività di
            progettazione architetturale (\textit{design}). Attivo dalla baseline RTB.
      \item \textbf{Programmatore:} Ruolo del gruppo responsabile delle attività di
            codifica del prodotto software; attivato durante la fase di sviluppo.
      \item \textbf{Project Manager:} Figura responsabile della pianificazione, esecuzione e chiusura di un progetto. Nel gruppo, questa funzione è ricoperta dal Responsabile.
      \item \textbf{Prompt:} Input testuale in linguaggio naturale fornito da un
            utente a un sistema di intelligenza artificiale per guidarne la
            generazione di contenuto.
      \item \textbf{Prompt Engineering:} Attività di progettazione e ottimizzazione dei prompt forniti a un LLM per ottenere risposte coerenti, accurate e aderenti agli obiettivi del sistema.
      \item \textbf{Proxy:} Design pattern strutturale che fornisce un surrogato o segnaposto per un altro oggetto, controllandone l'accesso e incapsulando dettagli come la comunicazione di rete. Nel progetto è utilizzato lato Frontend per invocare le funzionalità esposte dal Backend come se fossero locali.
      \item \textbf{Pull Model:} Variante del pattern Observer in cui l'osservatore, ricevuta la notifica di un cambiamento, interroga attivamente il Subject per ottenere lo stato aggiornato, anziché riceverlo direttamente nella notifica (Push Model). Adottata per l'architettura MVC del Frontend.
      \item \textbf{Pull Request:} Meccanismo di GitHub utilizzato per proporre
          l'integrazione di modifiche da un branch di lavoro al branch
          \texttt{main}. Nel flusso documentale del gruppo costituisce il punto
          di ingresso obbligatorio per la verifica: il Verificatore esamina le
          modifiche, lascia eventuali commenti e approva o rifiuta l'integrazione.
\end{itemize}
 
% ------------- Q -------------
\phantomsection
\addcontentsline{toc}{subsection}{Q}
\subsection*{Q}
\begin{itemize}[leftmargin=*]
      \item \textbf{Quality Metrics Compliance Rate (QMCR)}: Metrica di processo che misura
      la percentuale di metriche del cruscotto qualità che raggiungono il valore
      ottimale rispetto al totale delle metriche considerate.
\end{itemize}
 
% ------------- R -------------
\phantomsection
\addcontentsline{toc}{subsection}{R}
\subsection*{R}
\begin{itemize}[leftmargin=*]
      \item \textbf{Rate Limiting:} Tecnica che limita il numero di richieste effettuabili verso un servizio in un determinato intervallo temporale.
      \item \textbf{React:} Libreria JavaScript sviluppata da Meta per la creazione di interfacce utente reattive basate su componenti.
      \item \textbf{README:} File testuale che descrive informazioni operative relative
      a una cartella o componente del progetto, come prerequisiti, procedura di avvio,
      configurazione e problemi noti.
      \item \textbf{Redattore:} Ruolo del gruppo responsabile della produzione del
          contenuto di un documento, nel rispetto del template approvato e delle
          convenzioni redazionali definite nelle Norme di Progetto. Non può
          verificare il proprio lavoro.
      \item \textbf{Redo:} Funzionalità dell'editor che permette di ripristinare una modifica precedentemente annullata.      
      \item \textbf{Refactor:} Tipo di attività o commit che indica una modifica interna
      al codice o alla struttura del progetto senza cambiamento del comportamento
      osservabile.
      \item \textbf{Refactoring:} Attività di ristrutturazione del codice sorgente
      finalizzata a migliorarne la leggibilità, la manutenibilità e la qualità
      interna, senza modificarne il comportamento esterno. Utilizzata per ridurre
      il debito tecnico accumulato.
      \item \textbf{Render (o Rendering):} Processo di conversione e visualizzazione a schermo del codice sorgente (es. testo in Markdown) in un formato grafico leggibile e impaginato per l'utente finale.
      \item \textbf{Repository:} Archivio Git centralizzato che contiene documenti,
          codice sorgente e cronologia completa delle versioni del progetto.
          Ospitato su GitHub; il branch \texttt{main} contiene esclusivamente le
          versioni approvate e stabili.
      \item \textbf{Responsabile:} Ruolo del gruppo responsabile del coordinamento delle attività, dell'elaborazione di piani e scadenze, dell'approvazione del rilascio di prodotti parziali o finali e della rappresentanza del gruppo verso committente e proponente.
      \item \textbf{RESTful API:} Interfaccia di programmazione che segue i principi architetturali REST per la comunicazione tra client e server.
      \item \textbf{Review Coverage:} Metrica di processo che misura la percentuale di Pull Request approvate da un revisore diverso dall'autore.
      \item \textbf{Revisione:} Fase del ciclo di vita documentale in cui il
          Verificatore controlla struttura, coerenza, correttezza linguistica e
          conformità del documento alle Norme di Progetto, prima che esso venga
          pubblicato nel branch \texttt{main}.
      \item \textbf{Risorse Residue:} Ore e budget ancora disponibili per ciascun ruolo
      e per ciascun membro del team, calcolati sottraendo le ore effettivamente
      rendicontate da quelle preventivate. Monitorate al termine di ogni sprint
      tramite le tabelle di prospetto nel Piano di Progetto.
      \item \textbf{Routing:} Meccanismo che gestisce la navigazione e il collegamento tra diverse pagine o viste di un’applicazione.
      \item \textbf{Requirements Stability Index (RSI):} Indice di stabilità dei requisiti. 
      Misura la variazione dei requisiti nel tempo.
      \item \textbf{RTB (Requirements and Technology Baseline):} Prima revisione
          formale di avanzamento del progetto. Fissa i requisiti da soddisfare
          (Analisi dei Requisiti) e le tecnologie adottate, dimostrate tramite
          Proof of Concept. Si articola in due parti: valutazione del docente
          Cardin e presentazione al docente Vardanega.
      \item \textbf{Ruolo:} Funzione formalmente assegnata a un membro del gruppo
          per lo svolgimento di specifiche attività. I ruoli previsti sono:
          Responsabile, Amministratore, Analista, Progettista, Programmatore,
          Verificatore. Ogni membro deve ricoprire un solo ruolo alla volta e
          ruotare su tutti i ruoli nel corso del progetto.
\end{itemize}
 
% ------------- S -------------
\phantomsection
\addcontentsline{toc}{subsection}{S}
\subsection*{S}
\begin{itemize}[leftmargin=*]
      \item \textbf{SAL (Stato Avanzamento Lavori):} Momento di verifica periodica o documento che attesta lo stato di completamento delle attività di progetto rispetto alla pianificazione.
      \item \textbf{Sanificazione dei Dati:} Processo di rimozione o protezione di informazioni sensibili prima della loro elaborazione o trasmissione.
      \item \textbf{Scalabilità:} Capacità di un sistema software di adattarsi all’aumento del carico di lavoro mantenendo prestazioni adeguate.
      \item \textbf{Scenario Principale:} Sequenza di passi che descrive il flusso standard di un Caso d'Uso, dal punto di vista dell'attore.
      \item \textbf{Scenari Alternativi:} Varianti del flusso principale che descrivono eccezioni, estensioni o condizioni particolari del Caso d'Uso.
      \item \textbf{Schedule Variance (SV):} Varianza dei tempi.
      Valore positivo indica anticipo, negativo ritardo.
      \item \textbf{Scope:} Parte del progetto, documento o area tecnica a cui si riferisce
      un commit. Nel workflow del gruppo può indicare, ad esempio, un'area tecnica come
      \texttt{frontend} o un documento come \texttt{ndp}.
      \item \textbf{Scope Creep:} Fenomeno per cui l'ambito di un progetto si espande
      in modo incontrollato durante lo sviluppo, a causa dell'aggiunta di
      funzionalità non pianificate o non approvate formalmente. Costituisce uno
      dei rischi principali del progetto (RR2) e viene contrastato tramite un
      backlog rigoroso e l'approvazione obbligatoria del
      Responsabile per ogni nuova feature.
      \item \textbf{Scrum:} Modello di sviluppo agile basato su iterazioni brevi (sprint), ruoli definiti e revisione incrementale del prodotto. Nel progetto è adottato con rotazione settimanale dei ruoli.
      \item \textbf{Second Brain:} Nome del capitolato d'appalto scelto dal gruppo, proposto dall'azienda Zucchetti.
      \item \textbf{Sei Cappelli per Pensare:} Metodologia ideata da Edward de Bono che struttura il pensiero in sei prospettive distinte (logica, emotiva, critica, creativa, organizzativa e informativa). Nel progetto viene applicata tramite prompt dedicati per analizzare un testo secondo diverse angolazioni.
      \item \textbf{Semantic Versioning:} Convenzione di versionamento adottata dal
          gruppo basata sul formato \texttt{vX.Y.Z}, dove \texttt{X} (Major)
          identifica rilasci ufficiali o cambiamenti strutturali profondi,
          \texttt{Y} (Minor) indica aggiunte o modifiche significative di contenuto
          e \texttt{Z} (Patch) segnala correzioni ortografiche o modifiche non
          semantiche.
      \item \textbf{Simple Factory:} Pattern creazionale che centralizza in un'unica classe la logica di istanziazione di oggetti correlati, senza costituire un Factory Method GoF in senso stretto. Nel progetto è utilizzato per risolvere e istanziare l'operazione AI corretta a partire dal tipo richiesto, tramite un registro auto-registrante.
      \item \textbf{Single Point of Failure:} Componente o risorsa del progetto la cui
      indisponibilità causa il blocco dell'intera attività correlata. Nel contesto
      del gruppo si riferisce alla situazione in cui un task
      è assegnato a un unico membro senza documentazione condivisa, esponendo il
      team al rischio RO1.
      \item \textbf{Solo Editor:} Modalità che mostra esclusivamente l’area di editing.
      \item \textbf{Solo Render:} Modalità che mostra esclusivamente la versione renderizzata del documento.
      \item \textbf{Schedule Performance Index (SPI):} Indice di performance dei tempi.
      Se >1 si è in anticipo, se <1 si è in ritardo.
      \item \textbf{Spike Tecnico:} Attività esplorativa a tempo limitato svolta dal team
      per valutare la fattibilità o stimare la complessità di una tecnologia o di
      una funzionalità, prima di inserirla nel backlog come task
      pianificabile. Produce conoscenza, non deliverable.
      \item \textbf{Split View:} Modalità dell’editor che mostra simultaneamente la vista di editing e la vista renderizzata del documento.
      \item \textbf{Sprint:} Iterazione di lavoro breve e limitata nel tempo tipica delle metodologie agili, in cui il team si impegna a completare un set predefinito di task.
      \item \textbf{Sprint Backlog:} Insieme delle issue pianificate per lo sprint corrente, selezionate durante lo Sprint Planning e monitorate tramite GitHub Projects.
      \item \textbf{Sprint Planning:} Cerimonia agile svolta all'inizio di ogni
            sprint in cui il gruppo seleziona le attività da completare nel periodo
            e le assegna ai membri, definendo scadenze e priorità.
      \item \textbf{Sprint Retrospective:} Cerimonia agile svolta al termine di
            ogni sprint in cui il gruppo analizza il proprio modo di lavorare,
            identifica criticità e definisce azioni di miglioramento per lo sprint
            successivo.
      \item \textbf{Sprint Review:} Cerimonia agile svolta al termine di ogni sprint
            in cui il gruppo valuta i risultati ottenuti rispetto agli obiettivi
            pianificati, confrontando avanzamento atteso e avanzamento conseguito.
      \item \textbf{SQL:} Linguaggio standard utilizzato per la gestione e interrogazione di database relazionali.
      \item \textbf{Stack Tecnologico:} Insieme delle tecnologie, strumenti e framework utilizzati per sviluppare un sistema software.
      \item \textbf{Stakeholder:} Individuo, gruppo o organizzazione che ha un interesse attivo nel progetto e può influenzarne (o esserne influenzato da) i risultati, come il proponente o il committente.
      \item \textbf{State Management:} Insieme di tecniche e strumenti utilizzati per gestire lo stato applicativo in sistemi frontend complessi.
      \item \textbf{Strategy:} Design pattern comportamentale che definisce una famiglia di algoritmi intercambiabili, incapsulando ciascuno in una classe separata implementante un'interfaccia comune. Nel progetto è utilizzato per le operazioni AI (riassunto, traduzione, riscrittura, ecc.), ciascuna con la propria logica di costruzione del prompt.
\end{itemize}
 
% ------------- T -------------
\phantomsection
\addcontentsline{toc}{subsection}{T}
\subsection*{T}
\begin{itemize}[leftmargin=*]
      \item \textbf{T-y.z:} Codice identificativo delle slide del corso di Ingegneria del Software citate come riferimento informativo nei documenti di progetto. \texttt{y} è il numero progressivo della lezione e \texttt{z} l'eventuale numero di slide o sezione all'interno di essa.
      \item \textbf{Tailoring:} Processo di adattamento di uno standard (nel caso del
            gruppo: ISO/IEC 12207:1995) al contesto specifico del progetto, mediante
            la selezione motivata dei processi da attivare e l'esclusione giustificata
            di quelli non applicabili.
      \item \textbf{Task:} Unità operativa di lavoro assegnata a un membro del gruppo
      con ruolo, scadenza, stima del tempo previsto e tracciamento del tempo effettivo.
      Viene creata tramite Google Form, genera automaticamente una issue su GitHub ed
      è tracciata tramite GitHub Projects. A seconda della natura dell'attività, può
      essere documentale o non documentale e segue le convenzioni di naming definite
      nelle Norme di Progetto.
      \item \textbf{Task Completion Rate:} Metrica di processo che misura la percentuale
      di issue completate rispetto a quelle pianificate in uno sprint.
      \item \textbf{Technology Baseline:} Baseline che fissa le tecnologie, i
            framework e le librerie adottate per il progetto, dimostrando tramite
            Proof of Concept la loro adeguatezza e interoperabilità rispetto ai
            requisiti del capitolato. Costituisce uno degli obiettivi della
            revisione RTB.
      \item \textbf{Template:} Struttura predefinita e condivisa utilizzata per la
            redazione di tutti i documenti del gruppo, al fine di garantire uniformità
            formale e coerenza tra i diversi elaborati.
      \item \textbf{Template Method:} Design pattern comportamentale che definisce lo scheletro invariante di un algoritmo in una classe astratta, delegando alle sottoclassi la ridefinizione di singoli passi variabili. Nel progetto è utilizzato per l'esportazione della nota nei diversi formati (PDF, HTML, JSON).
      \item \textbf{Tempo Effettivo:} Tempo realmente impiegato da un membro del
            gruppo per completare un task, espresso in minuti. Registrato nel foglio
            di calcolo condiviso a consuntivo dell'attività.
      \item \textbf{Tempo Iniziale di Caricamento del Programma:} Metrica di prodotto
      che misura il tempo necessario affinché l'interfaccia principale dell'applicazione
      risulti disponibile e utilizzabile dall'utente.
            \item \textbf{Tempo Previsto:} Stima iniziale del tempo necessario per
            completare un task, espressa in minuti. Definita al momento della
            creazione del task e utilizzata come riferimento per il monitoraggio
            dell'avanzamento.
      \item \textbf{Ticket:} Segnalazione o richiesta di lavoro tecnica tracciata all'interno di un sistema di gestione. Nel flusso del gruppo, equivale a una Issue.
      \item \textbf{Time Efficiency:} Metrica di processo che misura la percentuale di
      ore produttive rispetto al totale delle ore spese.
      \item \textbf{Timeout LLM:} Errore generato quando il modello linguistico non risponde entro il tempo massimo previsto, gestito tramite messaggi di errore e retry.
      \item \textbf{Title Case:} Convenzione tipografica adottata dal gruppo per i
            titoli di sezioni, sottosezioni e paragrafi. Nella forma italianizzata adottata
            dal gruppo, le parole semanticamente rilevanti hanno iniziale maiuscola, mentre
            articoli, preposizioni e congiunzioni interne rimangono in minuscolo.
      \item \textbf{TLS:} Protocollo crittografico utilizzato per garantire sicurezza e riservatezza nelle comunicazioni di rete.
      \item \textbf{To Do:} Stato di una issue in GitHub Projects che indica che l'attività è stata pianificata ma non ancora iniziata.
      \item \textbf{TypeScript:} Superset di JavaScript che introduce tipizzazione statica e strumenti avanzati di sviluppo.
      \item \textbf{Token:} Elemento digitale utilizzato per autenticazione, autorizzazione o identificazione all’interno di un sistema.
      \item \textbf{Traduzione:} Funzionalità del sistema che consente di
            convertire il testo presente nell'editor in una lingua diversa
            tramite invocazione di un LLM. Le lingue supportate sono Inglese,
            Francese e Spagnolo come requisito obbligatorio; il Tedesco è
            previsto come estensione facoltativa. Sono escluse le lingue con
            alfabeti non latini.
      \item \textbf{Trunk-based Development:} Strategia di branching in cui il branch
      \texttt{main} rappresenta il ramo stabile principale e ogni modifica viene
      sviluppata su un branch di breve durata, poi integrato tramite Pull Request.
      Nel progetto non viene utilizzato un branch \texttt{develop}.
\end{itemize}
 
% ------------- U -------------
\phantomsection
\addcontentsline{toc}{subsection}{U}
\subsection*{U}
\begin{itemize}[leftmargin=*]
      \item \textbf{UML (Unified Modeling Language):} Linguaggio di modellazione
      standard utilizzato per rappresentare graficamente la struttura e il
      comportamento di un sistema software. Nel contesto del progetto è
      adottato per la rappresentazione dei diagrammi dei casi d'uso.
      \item \textbf{Undo:} Funzionalità dell'editor che permette di annullare l'ultima modifica effettuata.
      \item \textbf{Unione Discriminata (Discriminated Union):} Tipo dato algebrico, tipico di TypeScript, composto da più varianti mutuamente esclusive distinte da un campo discriminante, ciascuna contenitore di dati privo di comportamento polimorfico. Nel progetto è preferita al pattern State GoF per rappresentare lo stato di una richiesta AI, poiché la logica di transizione resta centralizzata in un'unica classe.
      \item \textbf{Usabilità:} Capacità di un sistema di essere utilizzato in modo efficace, efficiente e soddisfacente dall'utente finale.
      \item \textbf{Utenza Target:} Profilo dell'utente finale a cui il sistema
      è destinato, definito in accordo con il proponente. Nel contesto del
      progetto Second Brain, l'utente target è un utente giovane,
      tipicamente uno studente, che utilizza l'applicazione per prendere
      appunti e desidera essere supportato nella loro elaborazione tramite
      LLM.
\end{itemize}
 
% ------------- V -------------
\phantomsection
\addcontentsline{toc}{subsection}{V}
\subsection*{V}
\begin{itemize}[leftmargin=*]
      \item \textbf{Validazione:} Processo che accerta che quanto prodotto risponda alle esigenze della proponente e agli obiettivi del capitolato. In fase RTB riguarda soprattutto requisiti, casi d'uso, scelte tecnologiche e perimetro del Proof of Concept.
      \item \textbf{Valutazione Dei Capitolati:} Documento gestionale esterno prodotto
          in fase di candidatura che analizza e confronta i capitolati proposti
          dalle aziende proponenti, includendo il resoconto degli incontri
          esplorativi, i criteri di valutazione, i vantaggi e gli svantaggi di
          ciascuna proposta e la motivazione della scelta finale.
      \item \textbf{VarGroup:} Azienda proponente di uno dei capitolati d'appalto analizzati dal gruppo durante la fase di candidatura.
      \item \textbf{Variabili d’Ambiente:} Parametri di configurazione esterni al codice sorgente utilizzati per migliorare sicurezza e portabilità del software.
      \item \textbf{VE-y.x:} Codice identificativo univoco di una decisione adottata
          in un incontro esterno (con proponente o docente), registrata nel relativo
          verbale esterno. \texttt{y} è il numero progressivo del verbale esterno
          e \texttt{x} il numero progressivo della decisione al suo interno.
      \item \textbf{Velocity:} Misura della quantità di lavoro completata dal team in uno
      sprint, espressa in ore o story point. Viene monitorata sprint
      per sprint per rilevare andamenti anomali e per affinare le stime dei sprint
      successivi.
      \item \textbf{Verificatore:} Ruolo del gruppo responsabile del controllo di
          conformità di documenti e codice prodotti dagli altri membri. Deve essere
          sempre un membro diverso dal Redattore dell'elaborato in esame; non può
          verificare il proprio lavoro.
      \item \textbf{Virtual Machine:} Emulazione software completa di un sistema operativo eseguita su hardware virtualizzato.
      \item \textbf{VI-y.x:} Codice identificativo univoco di una decisione adottata
          in una riunione interna, registrata nel relativo verbale interno.
          \texttt{y} è il numero progressivo del verbale interno e \texttt{x} il
          numero progressivo della decisione al suo interno.
      \item \textbf{Vue.js:} Framework JavaScript progressivo utilizzato per sviluppare interfacce utente moderne, leggere e modulari.
\end{itemize}
 
% ------------- W -------------
\phantomsection
\addcontentsline{toc}{subsection}{W}
\subsection*{W}
\begin{itemize}[leftmargin=*]
    \item \textbf{Way of Working (WoW):} Insieme strutturato e documentato di
          regole, strumenti, procedure e metodologie adottate dal gruppo per
          organizzare e svolgere il lavoro in modo uniforme e tracciabile.
          Corrisponde operativamente alle Norme di Progetto e si evolve in modo
          incrementale secondo il principio just-in-time per tutta la durata del
          progetto.
    \item \textbf{WhatsApp:} Applicazione di messaggistica istantanea utilizzata
          dal gruppo come canale secondario per comunicazioni rapide, notifiche
          urgenti e convocazioni di riunioni straordinarie.
\end{itemize}
 
% ------------- X -------------
\phantomsection
\addcontentsline{toc}{subsection}{X}
\subsection*{X}
\textit{Nessun termine.}
 
% ------------- Y -------------
\phantomsection
\addcontentsline{toc}{subsection}{Y}
\subsection*{Y}
\textit{Nessun termine.}
 
% ------------- Z -------------
\phantomsection
\addcontentsline{toc}{subsection}{Z}
\subsection*{Z}
\begin{itemize}[leftmargin=*]
    \item \textbf{Zucchetti:} Azienda proponente del capitolato "Second Brain", per il quale il gruppo ha scelto di presentare la propria candidatura.
\end{itemize}